%-------------------------
% CV: Profil Umum - Bahasa Indonesia
% Author: Hafidz Mulia
% Compile: pdflatex -interaction=nonstopmode "CV_Hafidz_Mulia_General_ID_2026.tex"
%-------------------------
\documentclass[a4paper,11pt]{article}

%----------- PACKAGES -----------
\usepackage[utf8]{inputenc}
\usepackage[T1]{fontenc}
\usepackage{lmodern}
\usepackage[top=0.5in, bottom=0.5in, left=0.6in, right=0.6in]{geometry}
\usepackage{titlesec}
\usepackage{enumitem}
\usepackage[hidelinks]{hyperref}
\usepackage{xcolor}
\usepackage{fancyhdr}
\usepackage{tabularx}

%----------- COLORS -----------
\definecolor{accent}{HTML}{0F4C81}

%----------- PAGE STYLE -----------
\pagestyle{fancy}
\fancyhf{}
\renewcommand{\headrulewidth}{0pt}
\renewcommand{\footrulewidth}{0pt}
\fancyfoot[C]{\footnotesize\thepage}

%----------- SECTION FORMATTING -----------
\titleformat{\section}
  {\Large\bfseries\color{accent}}
  {}
  {0em}
  {}
  [\titlerule]

\titlespacing*{\section}{0pt}{13pt}{7pt}

%----------- CUSTOM COMMANDS -----------
\newcommand{\cvname}[1]{%
  {\centering\LARGE\bfseries #1\par\vspace{2pt}}%
}

\newcommand{\cvtitle}[1]{%
  {\centering\large\color{accent} #1\par\vspace{4pt}}%
}

\newcommand{\cvcontact}[1]{%
  {\centering\small #1\par\vspace{10pt}}%
}

\newcommand{\cventry}[4]{%
  \noindent\textbf{#1} \hfill \textit{#3} \\
  \textit{#2}
  \ifx&#4& \else \\ \small #4 \fi
  \vspace{10pt}
}

\newcommand{\cvproject}[5]{%
  \noindent\textbf{#1} \hfill \textit{#2} \\
  \textit{\small #3} \\
  {\small #4}
  \ifx&#5& \else \par\noindent{\small \href{#5}{GitHub}} \fi
  \vspace{10pt}
}

%----------- DOCUMENT START -----------
\begin{document}

%==================== IDENTITAS ====================
\cvname{Hafidz Mulia}
\cvtitle{Analisis Kebutuhan \textbar\ Pengembangan Sistem \textbar\ Koordinasi Tim}
\cvcontact{%
  \href{tel:+6285158335740}{+62 851 5833 5740} \quad\textbar\quad
  \href{mailto:hafidzmuliia@gmail.com}{hafidzmuliia@gmail.com} \quad\textbar\quad
  Surabaya, Indonesia \\
  \href{https://hafmul.site}{hafmul.site} \quad\textbar\quad
  \href{https://github.com/hafidzmulia-its}{github.com/hafidzmulia-its} \quad\textbar\quad
  \href{https://www.linkedin.com/in/hafidz-mulia/}{linkedin.com/in/hafidz-mulia}
}

%==================== RINGKASAN ====================
\section{Profil Singkat}
Lulusan Matematika ITS dengan pengalaman di pengembangan sistem, proyek digital, pengajaran, dan organisasi. Saya terbiasa mencari tahu masalah pengguna, menyusun kebutuhan, lalu membuat solusi yang dapat dipakai. Di luar pekerjaan teknis, saya juga pernah mengoordinasikan tim dan mendampingi orang belajar. Saya tertarik pada pekerjaan yang menggabungkan analisis masalah, perbaikan proses, dan kerja sama dengan orang lain.

%==================== PENGALAMAN ====================
\section{Pengalaman Kerja}

\cventry
  {Pengembangan Sistem (Penuh Waktu)}
  {PT Semen Indonesia Logistik (SILOG)}
  {Sep 2026 -- Sekarang}
  {Mengembangkan dan memelihara aplikasi web internal berbasis Laravel/Filament, memperbaiki bug pada sistem yang sudah digunakan, serta membantu tim dalam proses deployment dan pekerjaan terkait cloud/DevOps.}

\cventry
  {Pendiri \& Pengembang Aplikasi (Paruh Waktu)}
  {NechCode --- Agensi Solusi Digital}
  {Feb 2026 -- Sekarang}
  {Mendiskusikan masalah dan kebutuhan dengan klien di bidang pendidikan, layanan pendukung kesehatan, dan usaha kecil. Menyusun alur kerja dan fitur, lalu mengerjakan pengembangan, peluncuran, dan perbaikan aplikasi.}

\cventry
  {Pengajar Pemrograman (Paruh Waktu)}
  {Timedoor Academy}
  {Jan 2025 -- Sekarang}
  {Mengajar dasar pemrograman kepada anak dan remaja, menyesuaikan cara mengajar dengan kebutuhan siswa, dan membimbing mereka saat kesulitan memahami materi atau memperbaiki error.}

\cventry
  {Koordinator Asisten Laboratorium}
  {Departemen Matematika ITS (Pemrograman \& Komputasi Visual)}
  {Agu 2024 -- Des 2025}
  {Membantu mengatur kegiatan praktikum dan koordinasi asisten. Mendampingi lebih dari 100 mahasiswa saat mempelajari konsep pemrograman dan mengerjakan latihan.}

\cventry
  {Magang Pengembang Aplikasi Web}
  {SMA IT Al Uswah Surabaya}
  {Jan 2025 -- Feb 2025}
  {Mengumpulkan kebutuhan dari staf sekolah dan membuat sistem absensi untuk menggantikan pencatatan manual, mulai dari perancangan sampai uji coba bersama pengguna.}

%==================== PENDIDIKAN ====================
\section{Pendidikan}
\cventry
  {Sarjana Matematika}
  {Institut Teknologi Sepuluh Nopember (ITS), Surabaya}
  {2022 -- 2026}
  {IPK: 3,40 / 4,00. Skripsi tentang pengaturan partisi data berbasis memori untuk mempercepat pengambilan data deret waktu di PostgreSQL/TimescaleDB, dengan data lebih dari 50 juta baris.}

%==================== ORGANISASI ====================
\newpage
\section{Kepemimpinan \& Organisasi}

\cventry
  {Wakil Ketua I, Sentrum Katalis}
  {HIMATIKA ITS}
  {Feb 2025 -- Des 2025}
  {Membantu koordinasi dan pengelolaan organisasi beranggotakan lebih dari 150 orang. Ikut memperbaiki alur administrasi dan pelaporan internal menggunakan sistem digital.}

\cventry
  {Ketua Angkatan Matematika 2022}
  {Institut Teknologi Sepuluh Nopember}
  {Mei 2023 -- Sekarang}
  {Menjadi penghubung komunikasi dan membantu koordinasi kegiatan akademik maupun kebersamaan bagi lebih dari 150 teman seangkatan.}

\cventry
  {Kepala Departemen Kaderisasi}
  {LKKI Matematika ITS}
  {Jan 2024 -- Des 2024}
  {Menyusun kegiatan pengembangan anggota dan menjaga partisipasi serta keberlanjutan kegiatan organisasi.}

\cventry
  {Pemateri Workshop \& Pelatihan}
  {Berbagai Organisasi Kampus}
  {Jul 2024 -- Nov 2025}
  {Menyampaikan materi dasar pengembangan web, pengelolaan SDM, dan kepemimpinan kepada lebih dari 100 peserta dari beberapa kegiatan.}

%==================== PROYEK ====================
\section{Proyek Pilihan}

\cvproject
  {eRapor --- Sistem Pelaporan Organisasi}
  {2025}
  {Administrasi Organisasi \textbar\ Next.js, TypeScript, MongoDB}
  {Membuat dashboard sesuai peran pengguna dan laporan otomatis dalam format PDF/Excel agar pekerjaan pelaporan lebih rapi dan tidak perlu diulang secara manual.}
  {https://github.com/hafidzmulia-its/erapor}

\cvproject
  {Portal TPQ --- Sistem Pengelolaan Pendidikan}
  {2026}
  {Administrasi Pendidikan \textbar\ Laravel, MySQL}
  {Menggabungkan catatan perkembangan siswa, jadwal pengajar, dan pembuatan laporan dalam satu sistem untuk mengurangi pekerjaan administrasi.}
  {}

\cvproject
  {SIDEDIKK --- Platform Edukasi \& Skrining Kesehatan}
  {2026}
  {Aplikasi Ramah Pengguna \textbar\ Next.js, Laravel, PWA}
  {Membuat aplikasi yang mudah dipakai di ponsel agar pengguna bisa mengikuti alur skrining dan memahami informasi kesehatan dengan lebih jelas.}
  {}

%==================== KEMAMPUAN ====================
\section{Kemampuan}
\begin{tabularx}{\textwidth}{@{}l X@{}}
\textbf{Analisis \& Proyek:} & Menggali kebutuhan, menyusun alur kerja, memperbaiki proses, menyusun ruang lingkup pekerjaan, uji coba bersama pengguna \\[3pt]
\textbf{Tim \& Komunikasi:} & Koordinasi tim, mentoring, mengajar, memandu diskusi, komunikasi organisasi \\[3pt]
\textbf{Teknis \& Data:} & SQL, Python, PostgreSQL, MySQL, SQL Server, Next.js, Laravel, REST API, dashboard dan laporan \\[3pt]
\textbf{Bahasa:} & Bahasa Indonesia (aktif); Bahasa Inggris (membaca dan menulis materi kerja/teknis) \\
\end{tabularx}

\end{document}
